\section*{Análisis de la competencia}
\addcontentsline{toc}{section}{Análisis de la competencia}
Ahora vamos a utilizar las cinco fuerzas de Porter para poder analizar nuestra competencia, estas cinco fuerzas se usan para identificar y analizar las fuerzas que cambian el espíritu competitivo y la intensidad del mercado o industria. Permitiendo identificar la intensidad de la competencia y potenciales beneficios además de identificar factores clave que lo apartan de otros negocios \parencite{Investopedia02}.

\subsection*{Rivalidad entre competidores existentes}
\addcontentsline{toc}{subsection}{Rivalidad entre competidores existentes}
El mercado para aplicaciones financieras en Colombia es sumamente competitivo. Agencias como Nequi, Davivienda se enfocan en ofrecer servicios bancarios tradicionales con la ventaja de la digitalización mientras que Bconomy o Fintonic ofrecen servicios más orientados al control de gastos financieros y de salud financiera. También la competencia ha aumentado al surgimiento de otras propuestas locales e internacionales que buscan atraer la atención.

Para buscar destacar en este sector altamente competitivo, el plan de negocios busca hacer un enfoque especializado en el aprendizaje financiero estructurado a comparación de la gestión transaccional y el monitoreo, centrando más en ser una Edtech con aspectos de una Fintech dirigido especialmente a nómadas digitales, jóvenes entre 18 a 24 años y personas de bajos estratos. Lo que permitirá ofrecer de forma más específica y dedicada educación económica hacia estos sectores de la población a comparación de los servicios genéricos dados por los competidores.

\subsection*{Amenaza de nuevos entrantes}
\addcontentsline{toc}{subsection}{Amenaza de nuevos entrantes}
Los servicios financieros tienen barreras de entrada generalmente bajas en términos de infraestructura y tecnología, facilitando en cierta medida la entrada de nuevos competidores dentro del mercado. Esta accesibilidad puede incrementar la competencia y presión dentro del mercado.

Sin embargo, para nuevos entrantes, lograr una base sólida de clientes, junto a establecer una reputación confiable y el cumplimiento de normativas son factores críticos que representan una barrera para la entrada al mercado. El tiempo necesario para cumplir estos impedimentos, lograr credibilidad, cumplir con las normativas y lograr la experiencia para entregar servicios de alta calidad limitan la capacidad de entrada de nuevos jugadores en el corto plazo.

\subsection*{Poder de negociación de los proveedores}
\addcontentsline{toc}{subsection}{Poder de negociación de los proveedores}
Se depende de herramientas y plataforma de terceros tales como Google Ads y otro tipo de soluciones, además se otorgan a estas tecnologías cierto tipo de poder dentro de la negociación. El aumento de costos para las licencias y suscripciones pueden afectar en los costos operativos de la plataforma.

Adicionalmente, la competencia por el talento dentro de áreas como el desarrollo web y el diseño elevan costos y afectan a la calidad del servicio. Por ello, la plataforma buscará abordar el tema implementando estrategias que puedan atraer y retener profesionales talentosos, buscar alianzas con universidades y el sector público, y finalmente optimizar el uso de herramientas y de los recursos disponibles.

\subsection*{Poder de negociación de los clientes}
\addcontentsline{toc}{subsection}{Poder de negociación de los clientes}
Dentro del mercado, los clientes poseen una amplia gama de opciones, tales como las expuestas anteriormente, generando un alto poder de negociación, la capacidad de buscar y comparar experiencias y servicios de múltiples proveedores facilitan su capacidad de cambiar la plataforma en caso de que los resultados sean insuficientes.

Para desafiar esta barrera, nuestra plataforma se centrará en lograr obtener soluciones personalizadas a través de una personalización real impulsada por el machine learning, así como la entrega de resultados medibles en hábitos financieros. Asegurando un alto retorno en la inversión. Además, al adaptar la plataforma para guiar las especificaciones de ruta para cada cliente y ofrecer un servicio de gran calidad, se logrará fidelizar a los clientes y reducir su rotación.

\subsection*{Amenaza de productos sustitutos}
\addcontentsline{toc}{subsection}{Amenaza de productos sustitutos}
Una amenaza significativa para el proyecto es la proliferación de productos sustitutos, representados principalmente por el contenido educativo gratuito de finfluencers en redes sociales y los módulos de aprendizaje básicos integrados en aplicaciones bancarias tradicionales. Estos competidores ofrecen alternativas de costo cero que permiten a los usuarios gestionar su conocimiento de forma empírica y autónoma; no obstante, estas soluciones suelen ser genéricas, carecen de rigor técnico y no ofrecen la personalización algorítmica ni el soporte especializado que una solución de ingeniería dedicada puede proporcionar.

Nuestra plataforma mitiga este riesgo mediante una propuesta de valor centrada en la autoridad técnica y la adaptabilidad pedagógica. A diferencia de los consejos informales de las redes sociales o las herramientas rígidas de la banca, el ecosistema propuesto se diferencia por su capacidad de precisión predictiva, utilizando modelos de \textit{Machine Learning} para proyectar escenarios financieros reales que superan la naturaleza personal de los contenidos virales. Además, la plataforma ofrece una hipercontextualización diseñada específicamente para el marco legal y tributario de Bogotá, brindando una seguridad jurídica de la que carecen las aplicaciones globales.